\section{Exploración en RL profundo}

\subsection{De Bandit a MDP}

En un MDP, el problema de exploración es más complejo:
\begin{itemize}
    \item Las acciones no solo afectan la recompensa inmediata, sino también los estados futuros
    \item Es necesario explorar \textbf{pares estado--acción} y no solo acciones
    \item Algunos estados son difíciles de alcanzar en sí mismos y requieren una exploración dirigida durante mucho tiempo
\end{itemize}

\subsection{Exploración basada en conteo}

\begin{importantbox}{Count-Based Exploration}
Idea básica: otorgar una recompensa adicional de exploración a los pares estado--acción que se han visitado pocas veces:
$$r_{\text{explore}}(s, a) = \frac{\beta}{\sqrt{N(s, a)}}$$
En RL profundo, como el espacio de estados es continuo y de alta dimensión, no es posible contar directamente. Los métodos de aproximación más comunes incluyen:
\begin{itemize}
    \item Conteo con hash: mapear los estados a un espacio discreto mediante una función hash
    \item Modelo de densidad: entrenar un modelo de densidad $\hat{p}(s)$ y usar $-\log \hat{p}(s)$ como medida de novedad
    \item RND (Random Network Distillation): usar el error de predicción respecto de una red aleatoria como medida de novedad
\end{itemize}
\end{importantbox}

\subsection{Exploración basada en curiosidad}

Otro enfoque: recompensa de exploración = error de predicción del modelo. Si el error de predicción para un estado es grande, eso indica que el estado es «novedoso» y vale la pena seguir explorándolo.

\begin{warningbox}{La trampa de la curiosidad}
Los métodos de exploración basados en el error de predicción tienen un modo de falla conocido: el \textbf{problema de la televisión con ruido} (noisy TV problem). Si el entorno contiene aleatoriedad intrínsecamente impredecible (como la estática en la pantalla de un televisor), el error de predicción del modelo nunca disminuye, lo que hace que el agente quede «atraído» por esa aleatoriedad sin sentido y no pueda continuar con una exploración significativa.
\end{warningbox}

\subsection{Por qué los robots y los LLM rara vez exploran desde cero}

La conclusión que el curso da en esta parte es bastante pragmática: en MDP enormes como el control de robots y los LLM, realizar una exploración eficaz desde cero suele ser computacionalmente inasumible. En la industria y en la investigación de frontera, las estrategias más comunes son:
\begin{itemize}
    \item usar demonstrations o un base model preentrenado para reducir el espacio de búsqueda;
    \item proporcionar shaped rewards siempre que sea posible;
    \item reservar la exploración realmente difícil para subproblemas más pequeños y mejor estructurados.
\end{itemize}

\begin{figure}[H]
\centering
\includegraphics[width=0.9\textwidth]{slides-images/slide-13.jpg}
\caption{En robótica y en los LLM, normalmente se recurre al preentrenamiento, a las demostraciones y a shaped reward, en lugar de explorar desde cero}
\end{figure}

\subsection{Resumen de la sección}

La exploración en RL profundo requiere manejar medidas de novedad en espacios continuos de alta dimensión; los métodos count-based y los métodos basados en curiosidad tienen cada uno sus ventajas y limitaciones.
